\documentclass[Total.tex]{subfiles}

\begin{document}
	\chapter{Fundamental of Concepts in Graph Theory}
		
		\section{Graphs}
		
			\subsection{}
		
			\begin{Def}
				\begin{itemize}
					\item (Directed Graph)
					\item (Undirected Grpah)
				\end{itemize}
			\end{Def}
			
			\subsection{Vertexes}
			
			\begin{Def}
				内容...
			\end{Def}
			\subsection{Matrix Representation}
		\section{Path, Circle and Connectedness}
			\begin{Def}
				\begin{itemize}
					\item A walk in a graph $G$ is a sequence as following:
					\begin{gather*}
						内容...
					\end{gather*}
					where $e_i=(x_i,x_{i+1})$ or $e_{i}=\{x_i,x_{i+1}\}$
					\item A \textbf{path}
					\item A \textbf{cycle} in a graph $G$ is a 
				\end{itemize}
			\end{Def}
			\subsection{Connectedness}
			\subsubsection{Undirected Graph}
			\begin{lemma}
				Suppose $G$
			\end{lemma}
			\begin{proof}
				内容...
			\end{proof}
			
			\begin{Def}
				内容...
			\end{Def}
			\begin{proposition}
				内容...
			\end{proposition}
			\begin{proof}
				Take the negation $G$ is disconnected$\Leftrightarrow$ There exists some collection $X$ such that 
			\end{proof}
			
			Then, the component for graphs is well-defined
			\begin{Def}
				Let $G$ be a undirected graph, then,
				\begin{gather*}
					V(G)=\cup V(G_i)
				\end{gather*}
			\end{Def}
			\subsubsection{Directed Graph}
				\begin{proposition}
					内容...
				\end{proposition}
				\begin{proof}
					内容...
				\end{proof}
				\begin{Def}
					内容...
				\end{Def}
			\subsubsection{Euler Tour}
				\begin{Def}
					\begin{itemize}
						\item 
						\item 
					\end{itemize}
				\end{Def}
				\begin{proposition}
					A connected graph has one and only one Eulertour, when the graph is Euler.
				\end{proposition}
				\begin{proof}
					\begin{itemize}
						\item $\Rightarrow$
						\item $\Leftarrow$
					\end{itemize}
				\end{proof}
			\subsection{Leaf,Tree, Forest}
				\subsubsection{Tree, Leaves}
					\begin{Def} Suppose following graphs are undirected.
						\begin{itemize}
							\item (Tree) A
							\begin{itemize}
								\item Connected
								\item Undirected
								\item with no cycle
							\end{itemize}
							graph is called \textbf{tree}.
							\item The vertice of grad $1$ is called \textbf{leaf}.
						\end{itemize}
					\end{Def}
					\begin{lemma}
						The following proposition holds: For undirected graph $G$, we have
						\begin{itemize}
							\item Every tree with at least two vertices contains at least $2$ leaves;
							\item Suppose $G$ be a tree with 
							\begin{itemize}
								\item $n$ Vertices
								\item $m$ edges
								\item $p$ connected components.
							\end{itemize}
							Then we have
							\begin{gather*}
								n=m+p
							\end{gather*}
						\end{itemize}
					\end{lemma}
					\begin{proof}
						\begin{itemize}
							\item 
							\item 
						\end{itemize}
					\end{proof}
					
					\begin{proposition}
						The followings are equivalent:
						\begin{itemize}
							\item 
							\item 
							\item 
							\item 
							\item 
							\item 
							\item 
						\end{itemize}
					\end{proposition}
					\begin{proof}
						内容...
					\end{proof}
				\subsubsection{Arborenzence}
					\begin{Def}
						内容...
					\end{Def}
				\subsubsection{Minimal Spanned Tree}
					\begin{Def}
						内容...
					\end{Def}
			\subsection{Graph }
	\chapter{Optimal Tree and Path}
	
		\section{Minimum Spanning Tree Problem}
		
			\subsubsection{Kruskals Algorithms}
			
			\subsubsection{Prim Algorithms}
	
		\section{Graph Traversal}
		
		\begin{Def}
			\begin{itemize}
				\item A \textbf{connected and undirected graph}, which contains no circle, is called \textbf{tree}. 
				
				\item A \textbf{undirected graph}, which contains no circle, is called \textbf{forest};
				
				\item The vertex $e$, such that $|\delta(e)|=1$, is called \textbf{leaf},representing endpoints.
			\end{itemize}
		\end{Def}
		
		\begin{lemma}
			The following propositions hold:
			
			\begin{itemize}
				\item Every tree with at least $2$ vertexes contains 2 leaves;
				\item Suppose $G$ be a forest with $n$  vertexes, $m$ edges and $p$ components. Then: 
				
				\begin{gather*}
					n=m+p
				\end{gather*}
			\end{itemize}
		\end{lemma}
		
		\begin{proof}
			内容...
		\end{proof}
		
		\begin{proposition}
			Suppose $G$ be a undirected graph with $n$ vertexes. The following are equivalent
			
			\begin{itemize}
				\item $G$ is a tree;
				\item $G$ has $n-1$ vertexes and contains no circle;
				\item $G$ has $n-1$ vertexes and is connected;
				\item $G$ is connected and by deleting a vertex, we get a disconnected graph;
				\item $G$ is a graph with $\delta(X)\not=\varnothing$, for every $\varnothing\subseteq X\subseteq V(G)$, and the property is destroyed by deleting one vertex;
				\item $G$ is a forest. 
				\item Every two path in $G$ is connected by a unique path;
			\end{itemize}
		\end{proposition}
		
		\begin{proof}
			内容...
		\end{proof}
		
		\begin{Def}
			
			
			\begin{itemize}
				\item A directed graph $G$ is called \textbf{arborescence}, if:
				
				\begin{itemize}
					\item At every vertex $e\in E(G)$, the $\delta^+(e)\leq 1$;
					\item The correspondent undirected graph is a tree.
				\end{itemize}
				\item In an arborescence, there exists one and only one of vertex such that $\delta^-(r)=\varnothing$, which is called the \textbf{root};
				\item As a dual of concept forest, such directed graph is called \textbf{Silvanescence}.
			\end{itemize}
		\end{Def}
		
		\subsubsection{Graph Traversal Algorithm}
		
		\begin{proposition}
			内容...
		\end{proposition}	
		
		\begin{proof}
			内容...
		\end{proof}
		\subsubsection{DFS}
		\subsubsection{BFS}
		
		\subsubsection{Distance}
		
		\begin{Def}
			(Unweighted) For two vertexes $x,y$ in a graph $G$, the distance $dist(x,y)$ is the minimal number of $x-y$ path.
			
			If there exists no $x-y$ path, then $dist(x,y):=\infty$.
		\end{Def}
		\section{Weighted Graph}
		
		\begin{Def}
			A \textbf{weighted graph} is a graph $G$ with a function
			
			\begin{gather*}
				c:E(G)\rightarrow \mathbb{R}
			\end{gather*}
			
			which is called the \textbf{weight of edge}.
		\end{Def}
		
		Let $H\subseteq G$ be a subgraph. Then, 
		
		\begin{gather*}
			c(E(H)):=\sum_{e\in E(H)} c(e)
		\end{gather*}
		
		is called the \textbf{weight} of subgraph $H$.
		
		\subsubsection{Distance}
		
		
		
		A main problem is: \textbf{The shortest path problem}. 
		
		
		
		\subsubsection{Dijkstra's Algorithm}
		
		\begin{Def}
			\begin{itemize}
				\item 	Input: Weighted graph $(G,c)$ with $s\in V(G)$.  
				
				\item  Output: Arboriculture $A=(R,T)$ in $G$, such that $R$ contains all , and 
				
				\begin{gather*}
					disc_{(G,c)}(s,v)=disc_{(A,c)}(s,v);\forall v\in R
				\end{gather*}
			\end{itemize}
			
			The Pseudocode is given here: \ref{Dijkstras-Algorithm}
		\end{Def}
		
		\begin{algorithm}
			\caption{Dijkstra's Algorithm}
			\label{Dijkstras-Algorithm}
			\begin{algorithmic}
				\REQUIRE Weighted Graph $(G,c)$, Knot $s\in V(G)$;
				\ENSURE Arborescences $A=(R,T)$ in $G$, which contains all , and satisfy $dist_{G,c}(s,v)=dist_{(A,c)}(s,v)$
				\STATE R$\leftarrow$ 0;Q$\leftarrow\{s\}$ ;l(s)$\leftarrow$ 0\qquad$\backslash$$\backslash$ R:The ensured vertex with shortest path; Q:
				
				\WHILE{Q$\not=\varnothing$:}
				\STATE Select $v\in Q$ with $l(v)$ smallest;
				\STATE Q$\leftarrow Q-\{v\}$;R$\leftrightarrow R-\{v\}$;
				\STATE $\backslash\backslash$ To find out the neighbors of $v$;
				\FOR{$e=(v,w)\in \delta_G^+(v)$ with $w\notin R$} 
				\IF{$w\notin Q$ or $l(v)+c(e)<l(w)$}
				\STATE $l(w)\leftarrow l(v)+c(e);p(w)\leftarrow e;Q\leftarrow Q\cup \{w\}$;\qquad$\backslash$$\backslash$ If we find a shorter path from $s$ to $w$ via $v$,
				\STATE $\backslash$$\backslash$ then:$l(w):=l(v)+c(e);p(w)=e$, the previous edge; $Q:=Q\cup \{w\}$, extending the  
				\ENDIF
				\ENDFOR
				\ENDWHILE 
				
				$T\leftarrow \{p(v)|v\in R-\{s\}\}$;
			\end{algorithmic}
		\end{algorithm}
		
		
		\begin{proposition}
			Dijkstras Algorithmus is correct. And has a 
			
			\begin{gather*}
				\mathcal{O}(n^2+m)
			\end{gather*}
			
			where $n=|V(G)|,m=|E(G)|$.
		\end{proposition}
		
		\begin{proof}
			内容...
		\end{proof}
		
		\textbf{Remark}
		
		\textbf{Remark}
		
		\begin{Def}
			内容...
		\end{Def}
		
	
	\chapter{Flow in Network}
		\section{Network}
			\begin{Def}
				内容...
			\end{Def}
		
			
\end{document}